\section{A marked pressure comparison for the actual momentum fibers}
\label{sec:v99-pressure}
Normalize the initial lifted scatterer to the origin. Let
$\mathcal W_j(R)$ be the finite set of realizable physical itineraries
of $j$ flights, retaining the lifted displacement of each step. If
$[w]_R$ is the regular domain of $w$, define
\begin{equation}\label{eq:v99-word-sum}
 P_j(R)=\sum_{w\in\mathcal W_j(R)}
       \sup_{x\in[w]_R\cap M'}J^s_RT_R^j(x),\qquad P_0(R)=1.
\end{equation}
These are geometric upper sums, not replacements for the exact event.
They impose no occupation or section test; those tests are dropped
only on the upper side of the comparison below. There is one fixed
finite step alphabet of size $D$: successive lifted centers are at
distance at most $\tau_++2R_+$. It follows in particular that every
$P_j(R)$ is finite and that each fixed-$j$ supremum over $R$ is finite.

For comparison with the established pressure theory let
$\mathcal Q_j(R)$ be the stable-Jacobian sum at potential parameter
one over the homogeneous continuity partition
$\mathcal M_0^{j,H}$ of \cite[equations (2.4), (2.11)]{BDPressure}.
Its cells retain the true singularities through the terminal time and
homogeneity cuts through time $j-1$. Each cell determines one lifted
physical itinerary. Splitting into those cells can only increase a
sum of suprema of a positive function. Hence
\begin{equation}\label{eq:v99-pressure-domination}
                         P_j(R)\le\mathcal Q_j(R).
\end{equation}
No identification of the two sums is asserted.

Let $\mathcal K_m(R)$ denote \eqref{eq:v98-weighted-multiplicity}
with the radius held fixed. Thus $\mathcal K_m=\sup_R\mathcal K_m(R)$.

\begin{theorem}[Complete momentum-fiber pressure bound]
\label{thm:v99-pressure-comparison}
With a constant depending only on the uniform physical geometry,
\begin{equation}\label{eq:v99-pressure-comparison}
 \mathcal K_m(R)\le 1+C\left[ P_m(R)+D\sum_{i=1}^{m-1}q_0^{d_i}
                       P_i(R)P_{m-i-1}(R)\right].
\end{equation}
All physical preimages, internal incidences and original exact labels
are included. There is no factor counting homogeneity strips at the
marked collision.
\end{theorem}
\begin{proof}
Fix an almost-everywhere momentum target and one exact event. By
Theorem~\ref{thm:v99-unique}, each itinerary supplies at most one
preimage. The denominator in \eqref{eq:v98-gradient-density} is at
least $4\pi c_*\lambda_\circ^2$. The term without a mark is therefore
bounded by $CP_m(R)$ using Lemma~\ref{lem:v99-stable-weight}.

For the mark at $i$, use \eqref{eq:v99-remove-mark}. A full itinerary
is uniquely determined by its first $i$ flights, its single step at
time $i$, and its remaining $m-i-1$ flights, with the last block
translated back to the initial fundamental scatterer. The middle step
has at most $D$ choices. The weighted sum of its possible preimages
is consequently at most $CDP_i(R)P_{m-i-1}(R)$. This argument uses
physical itineraries, not the infinitely many homogeneity strips
whose inverse-incidence weight would no longer be summable after
removing the marked stable factor. Multiplying by $q_0^{d_i}$ and
summing proves the bound after taking the required essential and
exact-label suprema. The finite horizon alphabet and global
uniqueness supply every multiplicity used in this argument.
\end{proof}

\begin{corollary}[Bounded capacity at each fixed radius]
\label{cor:v99-fixed-radius}
For every fixed $R\in I$,
\begin{equation}\label{eq:v99-fixed-radius}
                    \sup_{m\ge1}\mathcal K_m(R)<\infty.
\end{equation}
In particular $m^6\mathcal K_m(R)\theta^{2m}\to0$ for every fixed
$0<\theta<1$.
\end{corollary}
\begin{proof}
For each fixed finite-horizon dispersing table, the pressure theorem
\cite[Lemma 3.8, proof]{BDPressure} gives
$\sup_j\mathcal Q_j(R)<\infty$. Apply
\eqref{eq:v99-pressure-domination} and
\eqref{eq:v99-pressure-comparison}, and use
$\sum_iq_0^{d_i}\le2q_0/(1-q_0)$. This use of the pressure theorem
is now justified by an explicit comparison of the same optical
weight and momentum preimages; it is not a comparison of cardinalities.
\end{proof}

The constant in the cited fixed-table pressure theorem has not been
shown here to be uniform as $R$ varies. The quantifiers in
\eqref{eq:v99-fixed-radius} must not be interchanged. The following
finite-block theorem gives an explicit family-uniform alternative.

\begin{theorem}[Uniform finite-block capacity certificate]
\label{thm:v99-block}
Put $\overline P_j=\sup_{R\in I}P_j(R)$. For every fixed integer
$L\ge1$ define
\[
 \rho_L=\max\{1,\overline P_L^{1/L}\},\qquad
 C_L=\max_{0\le j<L}\overline P_j\rho_L^{-j}.
\]
There is an explicit finite geometric multiple $A_L$ of
$1+C_L+DC_L^2\sum_{i\ge1}2q_0^i$ such that
\begin{equation}\label{eq:v99-block-growth}
                 \mathcal K_m\le A_L\rho_L^m\qquad(m\ge1).
\end{equation}
Consequently (SC) holds at a fixed auxiliary band whenever, for some
fixed $L$,
\begin{equation}\label{eq:v99-block-certificate}
                 \rho_L\theta_{B_*}^{\,2}<1.\tag{BC}
\end{equation}
If this certificate holds at an unbounded sequence of previously
fixed bands, all radius-uniform consequences of (SC) in the retained
article follow. No numerical verification of (BC) over $I$ is
claimed in this revision.
\end{theorem}
\begin{proof}
Splitting a physical itinerary into two consecutive blocks is
injective into the pair of block itineraries. The stable cocycle
gives $P_{j+k}(R)\le P_j(R)P_k(R)$, so also
$\overline P_{j+k}\le\overline P_j\overline P_k$.
Writing $m=vL+j$, $0\le j<L$, proves
$\overline P_m\le C_L\rho_L^m$. Insert this and the corresponding
bound for the two factors in \eqref{eq:v99-pressure-comparison}.
Since $\rho_L\ge1$ and the depth weights are summable, it gives
\eqref{eq:v99-block-growth}. Under (BC),
\[
 m^6\mathcal K_m\theta_{B_*}^{2m}
       \le A_Lm^6(\rho_L\theta_{B_*}^2)^m\longrightarrow0.
\]
All suprema over $R$, block lengths and bands are in the displayed
order. The block is fixed before $m$ and is allowed to depend on the
previously chosen auxiliary band, never on collision count.
\end{proof}

In particular
\[
 \limsup_m m^{-1}\log\mathcal K_m
       \le\inf_{L\ge1}\max\{0,L^{-1}\log\overline P_L\}.
\]
This is an actual uniform growth estimate expressed through geometric
partition sums. It replaces the unquantified fiber cardinality of
module 235. Establishing that its right side is zero, or certifying
the strict inequalities (BC), remains a family-uniform question;
pointwise-in-$R$ pressure bounds alone do not settle it.
